\documentclass[10pt]{article}
\usepackage[margin=20mm]{geometry}
\usepackage{amsmath,amssymb,mathtools,fontspec,unicode-math,xurl,hyperref}
\IfFontExistsTF{DejaVu Serif}{\setmainfont{DejaVu Serif}}{\setmainfont{Latin Modern Roman}}
\setmathfont{latinmodern-math.otf}
\setlength{\parindent}{0pt}\setlength{\parskip}{6pt}\sloppy
\title{Initial scientific review of
Bogomolov–Gieseker inequality for log terminal Kähler threefolds}
\author{Bounded initial review executor}
\date{29 October 2026}
\begin{document}
\maketitle
Abstract: The acquired arXiv v4 is reviewed in full, including its appendix. Two Major problems concern a false symmetry used in a central Green-kernel argument and an explicit metric construction which gives comparability without the asserted smooth orbifold extension. A separate Minor problem concerns the total mass of a regularized Monge–Ampère equation. Explicit calculations establish these defects. They do not establish a false Bogomolov–Gieseker theorem. A normalization repair is proved for the Green-kernel argument, and the dependence of the main theorem on a different interpolating metric is distinguished from the defective construction. This is a completed initial review, with scientific decision pending.\par
\section*{Source, scope, and principal conclusions}
The reviewed paper is by Henri Guenancia and Mihai Păun, with an appendix by Frédéric Campana, Andreas Höring, and Thomas Peternell [GP4]. The assigned identity is M25-0415; the journal DOI is https://doi.org/10.1002/cpa.70000. The exact acquired source is `arXiv:2405.10003v4`, submitted 2 March 2025, comprising 37 PDF pages. Its last page is dated 4 March 2025. Those two dates describe different source facts. Page numbers below refer to this PDF, whose printed pagination agrees with the physical pagination.\par
The principal theorem assumes a normal compact Kähler threefold with klt singularities, an ambient Kähler form \(\omega_X\), and a coherent reflexive \(\mathbb Q\)-sheaf \(\mathcal F\) of rank \(r>0\) which is stable with respect to \([\omega_X]\). It asserts \[\bigl(2r\,c_2(\mathcal F)-(r-1)c_1(\mathcal F)^2\bigr)\cdot[\omega_X]\geq0,\] with orbifold Chern numbers. The equality statement concerns local freeness and projective Hermitian flatness on \(X_{\mathrm{reg}}\), and extension of the projectivization after a finite quasi-étale Galois cover. It does not claim that the entire pulled-back reflexive sheaf becomes locally free.\par
All 37 pages were read. Definitions, parameter choices, hypotheses, theorem ranges, proofs, boundary terms, the appendix, and references were examined. All 37 source pages and 44 selected reference pages were also actually viewed as native rendered images. The three Problems below are proved defects of particular assertions or derivations. Neither a false principal theorem nor failure of every possible proof has been established.\par
\section*{Problem 1: the normalized Green kernel need not be symmetric}
Equation reference: eq:normal. Problem 1 (Major; invalid step in the principal analytic dependency).  On p. 6, the Green function is specified by \[
 \Delta_{\omega_\varepsilon}G_x^\varepsilon
 =\delta_x-V_\varepsilon^{-1},
 \qquad \sup_yG_x^\varepsilon(y)=-1.\] The proof of Proposition 3.3 then uses symmetry in the estimates immediately preceding (3.6) and (3.7), p. 8. The normalization in [eq:normal] (\#eq:normal) generally destroys that symmetry.\par
The operative setting is the regularization of a bounded canonical Monge–Ampère metric on a compact klt Kähler space. Proposition 3.3 also assumes Property (P): cutoffs whose \(L^{2+\epsilon}\) gradient norm and \(L^{1+\epsilon}\) Laplacian norm tend to zero, for a fixed \(\epsilon>0\). Its conclusion applies to nonnegative \(C^2\) functions for which the function and the negative part of its Laplacian belong to a sufficiently large \(L^p\) space. None of these conditions forces the maximum of a mean-zero Green kernel to be independent of its pole.\par
Equation reference: eq:K. Here is an exact example within the smooth compact Kähler case. Regard the unit round sphere as \(\mathbb P^1\), write \(z\) for its height, and choose \(0<\eta<1\). Put \[Y(z)=\frac{3z^2-1}{2},\qquad
 g_\eta=(1+\eta Y(z))g_0,\qquad H(z)=\frac{\eta Y(z)}{24\pi}.\] The density is positive, its integral is \(4\pi\), and \(\Delta_0Y=-6Y\), with the Laplacian convention \(\Delta G=\delta-(4\pi)^{-1}\). Define \[\begin{aligned}
 G_0(x,y)&=\frac{\log(1-x\cdot y)-\log2+1}{4\pi},\\
 K_\eta(x,y)&=G_0(x,y)+H(z_x)+H(z_y)-\frac{\eta^2}{120\pi}.
 \end{aligned}\] Since \(\Delta_{g_\eta}=(1+\eta Y)^{-1}\Delta_0\) in real dimension two, \(\Delta_yK_\eta=\delta_x^{g_\eta}-(4\pi)^{-1}\). Moreover, \(\int K_\eta(x,y)\,dA_{g_\eta}(y)=0\). For the latter identity, use \(\int G_0(x,y)Y(z_y)\,dA_0(y)=-Y(z_x)/6\) and \(\int Y^2\,dA_0=4\pi/5\). Thus [eq:K] (\#eq:K) is the symmetric mean-zero kernel.\par
Equation reference: eq:asym. Let \(m(x)=\max_y K_\eta(x,y)\), let \(P\) be a pole of the sphere, and let \(E\) be an equatorial point. For \(x=P\), both \(G_0(P,y)\) and \(H(z_y)\) attain their joint maximum at \(y=-P\). For \(x=E\), set \(t=E\cdot y\). The inequalities \[\log\frac{1-t}{2}\leq-\frac{1+t}{2},
 \qquad z_y^2\leq2(1+t)\] give \[K_\eta(E,y)-K_\eta(E,-E)
 \leq-\frac{(1-\eta)(1+t)}{8\pi}\leq0.\] Consequently \[m(P)-m(E)=2\bigl(H(1)-H(0)\bigr)=\frac{\eta}{8\pi}.\] The source normalization is precisely \(G_x(y)=K_\eta(x,y)-m(x)-1\), so that \[
 G_P(E)-G_E(P)=-\frac{\eta}{8\pi}\ne0.\] Take \(\omega_X\) to be the Kähler form of \(g_\eta\) and choose the smooth canonical-bundle metric whose associated measure is \(\omega_X\). The canonical potential is then zero; \(X\) is smooth and klt, and Property (P) holds with constant cutoffs. On the identity resolution one may take \(\widehat\omega=\omega_X\), obtaining the scaling family \((1+\varepsilon)\omega_X\). The same defect survives its limit. Thus the example does not rely on singular-space pathology.\par
\subsection*{Counterevidence, repair, and impact}
Self-adjointness gives symmetry for the mean-zero kernel. It does not give symmetry after pole-dependent shifts. Changing a fixed convention for the Laplacian changes the constants in the example but does not remove the dependence on the pole.\par
A repair of the relevant estimates is available. On each compact approximating manifold let \(K_\varepsilon\) be mean-zero and symmetric, and let \(m_\varepsilon(x)=\sup_yK_\varepsilon(x,y)\). Since \(m_\varepsilon\geq0\), the source’s uniform \(L^1\) bound and \(\inf_\varepsilon V_\varepsilon>0\) imply \[V_\varepsilon(m_\varepsilon(x)+1)
 =\int(-G_x^\varepsilon)\,\omega_\varepsilon^n\leq C.\] Choose one finite \(M\) exceeding every \(m_\varepsilon(x)+1\) and replace the kernel by \[\widetilde G_\varepsilon(x,y)=K_\varepsilon(x,y)-M.\] It is symmetric and at most \(-1\). The shift from \(G^\varepsilon\) is uniformly bounded, so the \(L^{1+\delta}\) estimate persists; its derivative in the second variable is unchanged. In addition, \[\frac{|\nabla_y\widetilde G_\varepsilon|^2}
      {(-\widetilde G_\varepsilon)^{1+\beta}}
 \leq
 \frac{|\nabla_yG_x^\varepsilon|^2}
      {(-G_x^\varepsilon)^{1+\beta}}.\] A common subsequence on compact subsets off the diagonal preserves symmetry in the limit. The cut-off estimates can then be performed with this kernel. The acquired Guo–Phong–Song–Sturm reference [Retained latex: \textbackslash{}cite[eq.\textasciitilde{}(5.17), p.\textasciitilde{}19]\{GPSS\}] uses a global infimum over both variables in its opposite-sign convention, exactly the appropriate kind of shift.\par
The false identity occurs in the proof of Proposition 3.3, which supplies Corollary 5.11, Theorem 5.15, and the subsequent current-closure argument for the main theorem. This makes the invalid derivation Major under the review criterion, even though the normalization repair is local and preserves the estimates. Equation [eq:asym] (\#eq:asym) is not a counterexample to Proposition 3.3 or to the final Chern-number inequality.\par
\section*{Problem 2: the explicit local potential is not orbifold smooth}
Problem 2 (Major; failure of the construction proving Theorem 4.3).  The definition on p. 9 requires a smooth Kähler metric after pullback to an orbifold chart. Theorem 4.3, pp. 11–12, asserts such a metric for a compact klt Kähler threefold whose orbifold singularities are ordinary double points. Claim 4.4 and (4.5), p. 12, prove comparability with the Euclidean metric. The construction and the three assertions on p. 13 do not supply the required smooth extension, and the claimed independence of a local embedding is false.\par
Equation reference: eq:levi. Let \[X=\mathbb P(1,1,2)\times\mathbb P^1.\] It is compact, projective, and klt. Every singularity is the product of an \(A_1\) surface quotient with a smooth curve; its non-orbifold locus \(Z\) is empty. Thus it satisfies the stated geometric assumptions. A local uniformization is \[\iota(u,v,z)=(u^2,v^2,uv,z),\qquad
 f(Z)=Z_1Z_2-Z_3^2.\] For \(c>0\) set \[F_c(u,v)=\sqrt{|u|^4+|v|^4+c|u|^2|v|^2}.\] Writing \(s=|u|^2\) and \(t=|v|^2\) and differentiating directly gives, for every real \(r>0\), \[
 \frac{\partial^2 F_c}{\partial u\,\partial\bar u}(r,0)=1,
 \qquad
 \frac{\partial^2 F_c}{\partial u\,\partial\bar u}(0,r)=\frac c2.\] Claim 4.4 uses \(c=1\), giving the distinct limits \(1\) and \(1/2\). The actual pullback of the gradient norm in (4.5) is \[|\nabla f|\circ\iota
 =\sqrt{|u|^4+|v|^4+4|u|^2|v|^2}=F_4,\] giving limits \(1\) and \(2\). Thus neither Hessian has a continuous extension at the origin. Euclidean comparability in Claim 4.4 can be true while smooth orbifold extension fails. The transverse coefficients of the pullback of a smooth ambient Kähler form are \(O(r^2)\) on these axes, so adding the prescribed large multiple of that form cannot remove the jump. A partition of unity with one chart equal to one near the point is an allowed local choice in the displayed construction and retains it.\par
There is also a direct counterexample to the embedding assertions. Embed the very same germ into \(\mathbb C^5\) by the ideal \((f,w)\), with \(w=0\). The prescribed sum of squared gradients is now \(1+|\nabla f|^2\), and the pulled-back potential is \[\widehat F(u,v)=\sqrt{1+|u|^4+|v|^4+4|u|^2|v|^2}.\] The original potential tends to zero and the new one tends to one. Hence their quotient has no uniform positive two-sided bound near the singular curve. Even subtracting the harmless constant one from \(\widehat F\) does not repair the asserted Hessian comparability, because \[\frac{\partial^2\widehat F}{\partial u\,\partial\bar u}(r,0)
 =\frac{2r^2+r^6}{(1+r^4)^{3/2}}
 =2r^2+O(r^6),\] whereas the corresponding coefficient of \(F_4\) equals one. The added ambient term is again \(O(r^2)\), leaving a degenerate transverse limit. These are two permitted embeddings of a germ satisfying the hypotheses, rather than an example from the excluded \(A_2\) case.\par
\subsection*{Counterevidence, repair, and impact}
For a suitable fixed chart, an actual orbifold potential is available: \[\sqrt{|Z_1|^2+|Z_2|^2+2|Z_3|^2}
 \quad\stackrel{\iota^*}{\longmapsto}\quad |u|^2+|v|^2.\] Adding \(|z|^2\) makes its pullback smooth and positive. This proves that the local existence assertion is not refuted. A replacement proof must nevertheless specify compatible potentials, justify their gluing, and verify the stated \(L^p\), \(p>1\), determinant condition on a resolution near the non-orbifold locus. Those global steps are not certified by the local replacement alone.\par
Theorem 4.3 is a substantive secondary result, and its explicit construction fails its claimed regularity and embedding invariance. The defect is therefore Major. It is not propagated indiscriminately to the Bogomolov–Gieseker theorem: Section 5 starts with a genuine orbifold metric away from a neighborhood of finitely many points and interpolates it with an ambient metric near those points. That construction is different from the gradient-norm formula in Theorem 4.3. The corresponding existence framework is also explicit in [Retained latex: \textbackslash{}cite[Theorem\textasciitilde{}1.4 and Section\textasciitilde{}6]\{CGN\}].\par
\section*{Problem 3: the smoothing changes the total mass}
Problem 3 (Minor; incorrect scalar normalization).  Equation (2.3), p. 5, has right-hand side \(c_\varepsilon\prod_i(|s_i|^2+\varepsilon^2)^{a_i}d\widehat V\), while p. 6 fixes \(c_\varepsilon=V_\varepsilon/V\). The hypotheses only imply \(\int\prod_i|s_i|^{2a_i}d\widehat V=V\). They do not imply the same identity after smoothing.\par
In the stated setting, \(\pi:\widehat X\to X\) is a log resolution of a compact klt Kähler space, \(a_i>-1\) are its discrepancies, \(|s_i|\) are smooth divisor norms, and \(V_\varepsilon=[\pi^*\omega_X+\varepsilon\widehat\omega]^n\). Integration of the equation requires \[c_\varepsilon I_\varepsilon=V_\varepsilon,\qquad
 I_\varepsilon=\int_{\widehat X}
       \prod_i(|s_i|^2+\varepsilon^2)^{a_i}d\widehat V.\]\par
A compact threefold example makes the incompatibility strict. Let \(X=\mathbb P(1,1,1,2)\). Its unique singularity is \(\mathbb C^3/\{\pm1\}\) of type \(\tfrac12(1,1,1)\). Blowing up its vertex gives a smooth resolution with exceptional divisor \(E\simeq\mathbb P^2\) and normal bundle \(\mathcal O_{\mathbb P^2}(-2)\). Its discrepancy is \(a=1/2\): adjunction gives \(K_{\widehat X}|_E=\mathcal O_{\mathbb P^2}(-1)\), whereas \(\pi^*K_X|_E\) is rationally trivial, so \(-2a=-1\). For the smooth positive volume form prescribed by the source, \[I_\varepsilon=\int_{\widehat X}\sqrt{|s|^2+\varepsilon^2}\,d\widehat V
 >\int_{\widehat X}|s|\,d\widehat V=V
 \qquad(\varepsilon>0).\] The printed right-hand side therefore has mass \(V_\varepsilon I_\varepsilon/V>V_\varepsilon\), while every left-hand side has mass \(V_\varepsilon\) by Stokes’ theorem. No smooth solution can satisfy that exact equation.\par
\subsection*{Counterevidence, repair, and impact}
Replace the scalar by \[c_\varepsilon=\frac{V_\varepsilon}{I_\varepsilon}.\] Since \(a_i>-1\) and the divisors have simple normal crossings, dominated convergence gives \(I_\varepsilon\to V\). The corrected factors tend to one and remain bounded above and away from zero. The uniform \(L^p\) estimate for a suitable \(p>1\) and smooth convergence away from the exceptional divisor are unaffected. This is an isolated normalization repair; the later approximation on p. 19 already describes its coefficient simply as a normalizing constant. The example refutes the printed regularized equation with the printed coefficient, not existence of the canonical metric or the final theorem.\par
\section*{Other operative assumptions and boundary checks}
The following checks delimit the findings and the scope of the review.\par
\subsection*{Conditional boundedness and approximation}
Property (P) remains a hypothesis for the canonical metric in Proposition B and Proposition 3.3 of v4. It is not silently supplied by the klt condition. The cutoff exponents in the proof can be met by taking \(p_0\) at least the maximum of the indicated conjugate exponents; in particular \(2p'\geq2p\) when \(p\) is conjugate to \(1+\epsilon\) and \(p'\) to \(1+\epsilon/2\). The finite-volume, nonnegative-function, and negative-Laplacian assumptions are essential to the stated application. The projectivity hypothesis of Theorem 4.1 is retained. Its use of Artin approximation, a finite cover, and an \(L^p\) density on a resolution was checked as a dependency argument, not as a new proof of every cited Monge–Ampère regularity theorem.\par
The hypotheses of the acquired Green-estimate reference include a cohomological intersection bound, entropy with exponent greater than the complex dimension, and a continuous lower density whose zero set is suitably small with connected complement [Retained latex: \textbackslash{}cite[pp.\textasciitilde{}2--3, 18--24]\{GPSS\}]. In the canonical resolution construction, the \(L^p\), \(p>1\), densities give the entropy bounds; the exceptional divisor has real codimension two. The bounds are compatible with the corrected mass normalization. The fixed positive volume lower bound was explicitly retained in Problem 1. The target’s references to later extensions of the Green estimates are not treated as new hypotheses absent from the acquired reference.\par
\subsection*{The interpolating metric and Hermitian–Einstein input}
For Section 5 the interpolation equals a fixed multiple of \(\omega_X\) near the finite set \(\{p_i\}\) and is genuinely orbifold outside a larger neighborhood. Cross terms are absorbed by the large ambient coefficient. The finite-cover/resolution Jacobian calculation in Proposition 5.6 gives integrability by choosing a sufficiently small positive power of the reciprocal Jacobian. The local uniformization cutoffs in Lemma 5.9 use fixed loci of complex codimension at least two; their exponents can be chosen small enough to give both required norms.\par
The Hermitian–Einstein existence and curvature estimates in Theorem 5.12 were compared with the actual Theorems 1.2 and 1.4 of [CGN]. The reference metric on a reflexive sheaf is a quotient metric, and its pullback is not automatically uniformly equivalent to a smooth bundle metric at the branching set. The power bounds and subsequent boundedness argument in Theorem 5.15 address this distinction. There is no valid general inference from a one-sided comparison of bundle metrics to a two-sided comparison of all endomorphism norms. For curvature, however, an algebraic check prevents a spurious objection: on real tangent vectors it is skew-Hermitian for the Chern metric, so its intrinsic squared norm is the sum of squared moduli of its eigenvalues, bounded by its Hilbert–Schmidt norm for any other metric. Thus the quoted curvature bound is not refuted by arbitrary non-normal endomorphism examples.\par
\subsection*{Local products, current closure, and equality}
Graf–Kirschner’s Lemma 5.8 gives a surface-quotient product after removal of a codimension-three analytic set [GK, p. 896]; in dimension three this set is finite. It should be included in \(\{p_i\}\). The surface quotient need not be a hypersurface in \(\mathbb C^3\). For example, \(\mathbb C^2/\mu_3\) with scalar action has four independent degree-three generators \(u^3,u^2v,uv^2,v^3\). Accordingly the three displayed transverse coordinate functions on p. 26 must be interpreted with an arbitrary finite number of generators in the general klt case. The subsequent Lemma 5.19 is intrinsically two-dimensional and its logarithmic integral argument does not require three ambient generators. The same derivative bounds follow for a finite sum of their squared norms. This is a clarification of the local model; no failure of the intrinsic integral estimate was proved.\par
The logarithmic cutoff integrals were independently checked using simultaneous principalization of the finitely many analytic ideals. After resolution, the required radial factors have the form \(dr/(r|\log r|)\) on logarithmic annuli; their integrals are uniformly bounded when the endpoints have fixed ratio in the logarithmic variable. For the final boundary integral at \(p_i\), the metric equals an ambient metric, the potential has gradient of order \(|Z|\), and curvature is \(L^2\). These facts give a bounded cutoff/potential factor and an integral over shrinking sets tending to zero by absolute integrability. The exterior derivative in Theorem 5.17 is \(\bar\partial\), as confirmed visually; treating it as \(\partial\) would create an irrelevant bidegree objection.\par
The equality argument uses nonnegative trace-free curvature density. Its reflexive extension across smooth isolated points and its projective-bundle conclusion are consistent with that range. The non-Cartier line-sheaf example in Remark 5.5 illustrates why the projectivization statement is weaker than local freeness of the whole pulled-back sheaf. It is not a stable-sheaf counterexample to the main theorem.\par
\subsection*{Appendix and Chern-number conventions}
The appendix’s abundance theorems retain normality, \(\mathbb Q\)-factoriality, terminal singularities, and nefness of \(K_X\); the positivity step also retains the stated numerical-dimension assumption. Proposition A.5 and Lemma A.6 concern the non-algebraic case. The lemma explicitly excludes projective uniruled klt examples with nonnegative Kodaira dimension, so those examples do not refute it. The resolution argument pairs exceptional divisors with the square of a pulled-back nef class and uses the fact that their images have codimension at least two.\par
For numerical Chern classes the precise general construction in [Retained latex: \textbackslash{}cite[Definition\textasciitilde{}5.2 and Section\textasciitilde{}5.2]\{GK\}] is a functional on \(H^{2n-4}(X,\mathbb R)\), obtained from compactly supported cohomology of the orbifold locus. It is not a general theorem identifying that functional with a unique class in \(H^4(X,\mathbb R)\). The integral definition in the target’s (5.3)–(5.4) agrees with the functional construction. This review uses that definition and does not claim a general \(H^4\)-extension theorem from the abbreviated wording in the appendix. No counterexample satisfying the appendix’s full preceding \(\mathbb Q\)-factorial setup was established.\par
\section*{Supporting checks}
The current arXiv record [record] explicitly says that an earlier result was removed because of a proof gap while the main result was unchanged. The metadata does not name a theorem number. A direct comparison of the acquired v3 introduction [Retained latex: \textbackslash{}cite[pp.\textasciitilde{}1--4]\{GP3\}] with v4 shows that v3 stated general Kähler orbifold regularity as Theorem B and a local Property (P) assertion in Proposition C.2; v4 keeps the relevant boundedness assertion conditional and identifies the general Kähler analogue of Theorem 4.1 as unknown in Remark 4.2. The withdrawn assertion is not counted as a current defect. This is a comparison of the specified introductions, not a full scientific review of v3 or a line-by-line comparison with the journal version.\par
`277016e83ae288ceb4dbe01b0d2a806dbbd5f70b0570a6897eb19b479152d728`.\par
The source-manifest SHA256 is\par
`0857f9fb4835263498fe2ebb55ab8610aeab2315723bd35bbf8be7b7d13695f8`.\par
The independent Python check uses exact rational arithmetic for algebraic constants and ordinary floating-point evaluations for corroboration. For \(\eta=1/2\), [eq:asym] (\#eq:asym) gives approximately \(-0.01989436788648692\). The grid evaluations do not prove the maximum identity; the inequalities above do. The directional Levi coefficients and the redundant-embedding degeneration are also established analytically. The Monge–Ampère mass discrepancy is an analytic integration argument, not a claimed numerical integration on a weighted projective variety. The saved result SHA256 is\par
`f375f4bbf7370f758eeec3651f92e8ad2e2ffe9e2a63f9640b938e847687640e`.\par
No author code was executed. No absence of code is used as a mathematical criticism. Public metadata and the relevant cited primary sources were checked; unsuccessful publisher/landing-page requests remain in the provenance record. No universal assertion that all possible corrections or supplements are absent is made.\par
The principal executor performed the scientific reading, calculations, and image inspections. The helper performed only acquisition, file binding, extraction, and rendering. The runtime self-report and the controller’s product-interface observation are recorded separately; neither is presented as independent backend certification. The final compilation, unchanged dependency checks, PDF parsing, page renders, and actual inspections are bound in the accompanying provenance and support inventory.\par
The disposition is two Major problems and one Minor problem, with the repair and impact boundaries stated above. There is no demonstrated false final Bogomolov–Gieseker inequality. Scientific decision remains pending, human reviews are empty, and this initial report is not independent acceptance.\par
9 H. Guenancia and M. Păun, with an appendix by F. Campana, A. Höring, and T. Peternell, Bogomolov–Gieseker inequality for log terminal Kähler threefolds, arXiv:2405.10003v4 (2 March 2025), 37 pp. Journal identifier: DOI 10.1002/cpa.70000. B. Guo, D. H. Phong, J. Song, and J. Sturm, Diameter estimates in Kähler geometry, arXiv:2209.09428v1 (20 September 2022), 50 pp. Examined pp. 1–3 and 18–25. The target cites the 2024 journal article; no equality of its text with this acquired preprint is asserted. J. Cao, P. Graf, P. Naumann, M. Păun, T. Peternell, and X. Wu, Hermite–Einstein metrics in singular settings, arXiv:2303.08773v2 (21 January 2024), 44 pp. Examined pp. 1–6 and 20–34. P. Graf and T. Kirschner, Finite quotients of three-dimensional complex tori, Ann. Inst. Fourier 70 (2020), 881–914. DOI 10.5802/aif.3326. Published PDF, including cover; examined physical pp. 13–20 (printed pp. 892–899). H. Guenancia and M. Păun et al., Bogomolov–Gieseker inequality for log terminal Kähler threefolds, arXiv:2405.10003v3 (18 July 2024), 43 pp. Only pp. 1–4 examined for version boundaries. Authors’ arXiv record for 2405.10003, https://arxiv.org/abs/2405.10003, accessed 29 October 2026. The version-specific PDFs and saved retrieval records identify the materials actually used.\par
\end{document}
